% checks_d2.tex -- d = 2 paragraph of the independent-checks section (from HT2dPert, 2026-10-14). Checks p2.*.
\paragraph{Two dimensions.}\label{htp:d2:checks}
Two solvers were written from a first-principles specification of the $d=2$ model (Sec.~\ref{htp:d2:model}) by agents that saw
neither the production code nor each other's work. The first is plain Rayleigh--Schr\"odinger perturbation theory to $\eps^8$
with full wavefunction corrections in exact rational arithmetic, in two codes (\texttt{fractions} over the full space, and a
reflection-reduced code modulo many primes with rational reconstruction checked on held-out primes), for
$L=0,\dots,\NTwoRSLmax$ in all three schemes. The second uses no perturbation theory: the truncated Hamiltonian is built at
complex $\eps$ on a circle with 256-bit arithmetic, $r$ and $g$ are Newton-solved, the eigenvalues are followed from the
references, and Cauchy integrals give the coefficients, for $L\le\NTwoCauchyLmax$. Over $C_n$, $r_n$, $g_n$, the reference
energies and all sub-cutoff vacua, $n\le8$, the production solver agrees with the exact one to $\HTpTwoImplAgree$ (scaled
difference $\max|a-b|/\max(1,|b|)$; raw and consistent $L\le\NTwoRSLmax$, literal $L\le14$), the Cauchy solver with the exact one
to $\NTwoCauchyAgree$, and the production $C_n$ with the Cauchy ones to $\NTwoPkgCauchy$. The large-$L$ sum codes of
Sec.~\ref{htp:d2:solvers} reproduce these values at small $L$. The raw scheme was specified from first principles; for the
consistent and literal schemes the definitions were taken from the production package, so there the check covers the
implementation, not the choice of window. These checks verify that we compute the stated truncation; they do not test its
continuum limit.
